\documentclass{article}
\usepackage{amsmath,amssymb,amsthm}
\usepackage{algorithm}
\usepackage{algpseudocode}
\usepackage{enumitem}
\usepackage{hyperref}
\newtheorem{theorem}{Theorem}
\newtheorem{lemma}{Lemma}
\newtheorem{assumption}{Assumption}
\newcommand{\norm}[1]{\left\lVert#1\right\rVert}
\newcommand{\R}{\mathbb{R}}

\begin{document}

\section*{Algorithm Name}
\textbf{Trust‑Region Gradient Method (TR‑GM)} for convex and $L$‑smooth objectives.

---------------------------------------------------------------------

\section*{Mathematical Formulation}
Let $f:\R^{n}\to\R$ be convex and continuously differentiable with $L$‑Lipschitz gradient:
\[
\norm{\nabla f(x)-\nabla f(y)}\le L\norm{x-y},\qquad\forall x,y\in\R^{n}. \tag{A1}
\]

At iteration $k$ we are given the current iterate $x_{k}$ and a trust‑region radius $\Delta_{k}>0$.
We consider the quadratic model
\[
m_{k}(p)=f(x_{k})+\nabla f(x_{k})^{\top}p+\frac{L}{2}\norm{p}^{2},
\]
and (approximately) solve the subproblem
\[
\min_{p\in\R^{n}}\; m_{k}(p)\quad\text{s.t.}\quad \norm{p}\le \Delta_{k}. \tag{TR}
\]

A closed‑form approximate solution is obtained by a scaled gradient step:
\[
p_{k}= 
\begin{cases}
-\displaystyle\frac{\Delta_{k}}{\norm{\nabla f(x_{k})}}\;\nabla f(x_{k}), & 
\text{if }\norm{\nabla f(x_{k})}>\Delta_{k},\\[2ex]
-\nabla f(x_{k}), & \text{otherwise}.
\end{cases} \tag{1}
\]

Define the *predicted reduction*
\[
\operatorname{pred}_{k}=m_{k}(0)-m_{k}(p_{k})
 = -\nabla f(x_{k})^{\top}p_{k}-\frac{L}{2}\norm{p_{k}}^{2}, \tag{2}
\]
and the *actual reduction*
\[
\operatorname{ared}_{k}=f(x_{k})-f(x_{k}+p_{k}). \tag{3}
\]

Given a parameter $\eta\in(0,1)$, we accept the step if
\[
\operatorname{ared}_{k}\ge \eta\,\operatorname{pred}_{k}. \tag{4}
\]
When the step is accepted we set $x_{k+1}=x_{k}+p_{k}$ and possibly enlarge the radius,
otherwise we keep $x_{k+1}=x_{k}$ and shrink the radius.

---------------------------------------------------------------------

\section*{Pseudocode}
\begin{algorithm}[H]
\caption{Trust‑Region Gradient Method (TR‑GM)}
\begin{algorithmic}[1]
\Require Initial point $x_{0}\in\R^{n}$, initial radius $\Delta_{0}>0$, 
parameters $\eta\in(0,1)$, $\gamma_{\text{inc}}>1$, $\gamma_{\text{dec}}\in(0,1)$,
maximum radius $\Delta_{\max}>0$.
\State $k\gets0$.
\While{stopping criterion not met}
    \State Compute gradient $g_{k}\gets\nabla f(x_{k})$.
    \If{$\norm{g_{k}}>\Delta_{k}$}
        \State $p_{k}\gets -\dfrac{\Delta_{k}}{\norm{g_{k}}}\,g_{k}$ \Comment{scaled step}
    \Else
        \State $p_{k}\gets -g_{k}$.
    \EndIf
    \State Compute predicted reduction $\operatorname{pred}_{k}$ by \eqref{2}.
    \State Compute actual reduction $\operatorname{ared}_{k}=f(x_{k})-f(x_{k}+p_{k})$.
    \If{$\operatorname{ared}_{k}\ge \eta\,\operatorname{pred}_{k}$}
        \State $x_{k+1}\gets x_{k}+p_{k}$ \Comment{accept}
        \State $\Delta_{k+1}\gets\min\{\gamma_{\text{inc}}\Delta_{k},\Delta_{\max}\}$.
    \Else
        \State $x_{k+1}\gets x_{k}$ \Comment{reject}
        \State $\Delta_{k+1}\gets\gamma_{\text{dec}}\Delta_{k}$.
    \EndIf
    \State $k\gets k+1$.
\EndWhile
\State \Return $x_{k}$.
\end{algorithmic}
\end{algorithm}

---------------------------------------------------------------------

\section*{Dry Run Example}
Consider the one‑dimensional quadratic
\[
f(w)=(w-3)^{2},\qquad \nabla f(w)=2(w-3).
\]
Choose $w_{0}=0$, initial radius $\Delta_{0}=1$, $\eta=0.1$, $\gamma_{\text{inc}}=2$, $\gamma_{\text{dec}}=0.5$.
Since $f$ is $L$‑smooth with $L=2$, we use the above algorithm.

\begin{center}
\begin{tabular}{c|c|c|c|c}
\textbf{Iter} & $w_{k}$ & $g_{k}= \nabla f(w_{k})$ & $p_{k}$ (from (1)) & $f(w_{k})$\\\hline
0 & $0$ & $-6$ & $p_{0}= -\dfrac{1}{6}(-6)= -1$ & $9$\\
  & & & $\operatorname{pred}_{0}= -g_{0}p_{0}-\frac{L}{2}p_{0}^{2}=6\cdot1-\frac{2}{2}\cdot1=5$ & \\
  & & & $\operatorname{ared}_{0}=f(0)-f(-1)=9-16= -7$ (reject) & \\
  & & & $\Delta_{1}=0.5$ & \\
\hline
1 & $0$ & $-6$ & $p_{1}= -\dfrac{0.5}{6}(-6)= -0.5$ & $9$\\
  & & & $\operatorname{pred}_{1}=6\cdot0.5-\frac{2}{2}\cdot0.25=3-0.25=2.75$ & \\
  & & & $\operatorname{ared}_{1}=f(0)-f(-0.5)=9-12.25=-3.25$ (reject) & \\
  & & & $\Delta_{2}=0.25$ & \\
\hline
2 & $0$ & $-6$ & $p_{2}= -\dfrac{0.25}{6}(-6)= -0.25$ & $9$\\
  & & & $\operatorname{pred}_{2}=6\cdot0.25-\frac{2}{2}\cdot0.0625=1.5-0.0625=1.4375$ & \\
  & & & $\operatorname{ared}_{2}=f(0)-f(-0.25)=9-10.5625=-1.5625$ (reject) & \\
  & & & $\Delta_{3}=0.125$ & \\
\end{tabular}
\end{center}
The example shows that, because the initial point is far from the minimiser, the trust‑region radius is repeatedly reduced until the step becomes acceptable. After a few more iterations the radius will be small enough that the step is taken and the algorithm converges to $w^{*}=3$.

---------------------------------------------------------------------

\section*{Assumptions}
\begin{assumption}[Convexity and Smoothness]\label{as:convex}
$f:\R^{n}\rightarrow\R$ is convex and differentiable, and its gradient is $L$‑Lipschitz:
\[
f(y)\le f(x)+\nabla f(x)^{\top}(y-x)+\frac{L}{2}\norm{y-x}^{2},\qquad\forall x,y\in\R^{n}. \tag{A1}
\]
\end{assumption}
\begin{assumption}[Bounded Level Set]\label{as:bounded}
For the initial point $x_{0}$ the level set 
\[
\mathcal{L}:=\{x\in\R^{n}\mid f(x)\le f(x_{0})\}
\]
is bounded, i.e., there exists $D>0$ such that $\norm{x-x^{*}}\le D$ for all $x\in\mathcal{L}$, where $x^{*}$ denotes any global minimiser.
\end{assumption}
\begin{assumption}[Algorithmic Parameters]\label{as:param}
\[
\eta\in(0,1),\qquad
\gamma_{\text{inc}}>1,\qquad
0<\gamma_{\text{dec}}<1,\qquad
0<\Delta_{\min}\le\Delta_{0}\le\Delta_{\max}<\infty .
\]
\end{assumption}

---------------------------------------------------------------------

\section*{Main Convergence Theorem}
\begin{theorem}[Global $O(1/k)$ rate]\label{thm:rate}
Assume \ref{as:convex}–\ref{as:param}. Let $\{x_{k}\}$ be the sequence generated by TR‑GM.
Define $f^{*}= \inf_{x}f(x)$. Then for all $K\ge1$
\[
f(x_{K})-f^{*}\;\le\;
\frac{L D^{2}}{2\,\eta\,K},
\]
where $D$ is the bound from Assumption~\ref{as:bounded}.
Consequently $\displaystyle\lim_{k\to\infty}f(x_{k})=f^{*}$.
\end{theorem}

---------------------------------------------------------------------

\section*{Theoretical Properties}
\begin{enumerate}[label=\arabic*.]
\item \textbf{Stability.}  Because the step $p_{k}$ always satisfies $\norm{p_{k}}\le\Delta_{k}$ and $\Delta_{k}\le\Delta_{\max}$, all iterates remain inside the bounded level set $\mathcal{L}$, guaranteeing numerical stability.
\item \textbf{Hyper‑parameter Sensitivity.}
    \begin{itemize}
        \item Larger $\eta$ makes the acceptance test stricter, potentially leading to more radius reductions but also to higher quality steps.
        \item $\gamma_{\text{inc}}$ controls how fast the radius can grow after a successful step; too large a value may cause oscillations.
        \item $\gamma_{\text{dec}}$ governs shrinkage; values close to $1$ make the algorithm conservative, while very small values may stall progress.
    \end{itemize}
    The $O(1/k)$ bound holds for any choice within the admissible ranges of Assumption~\ref{as:param}.
\item \textbf{Comparison to Gradient Descent.}
    If we fix $\Delta_{k}\equiv\infty$, the subproblem reduces to the exact gradient step $p_{k}=-\nabla f(x_{k})$, and the algorithm becomes standard gradient descent with step size $1/L$, achieving the same $O(1/k)$ rate. The trust‑region mechanism automatically adapts the step size and can be more robust when $L$ is unknown.
\end{enumerate}

---------------------------------------------------------------------

\section*{Discussion}
The theorem shows that for convex $L$‑smooth objectives the trust‑region method enjoys the same worst‑case convergence rate as optimal first‑order methods, while offering the practical advantage of an adaptive step‑size mechanism. In non‑convex settings the same analysis yields a bound on the norm of the gradient, but that is beyond the present scope.

---------------------------------------------------------------------

\section*{Preliminaries (Definitions and Lemmas)}
\begin{lemma}[Descent Lemma]\label{lem:descent}
If $f$ satisfies (A1) then for any $x\in\R^{n}$ and any $p\in\R^{n}$,
\[
f(x+p)\le f(x)+\nabla f(x)^{\top}p+\frac{L}{2}\norm{p}^{2}. \tag{L1}
\]
\end{lemma}
\begin{proof}
Directly from the definition of $L$‑Lipschitz gradient (A1). \qedhere
\end{proof}

\begin{lemma}[Predicted reduction lower bound]\label{lem:pred}
For the step $p_{k}$ defined in \eqref{1} we have
\[
\operatorname{pred}_{k}\;\ge\;
\frac{1}{2}\min\!\Big\{\,\norm{\nabla f(x_{k})}^{2},\,L\Delta_{k}^{2}\Big\}. \tag{L2}
\]
\end{lemma}
\begin{proof}
Two cases:
\begin{itemize}
    \item If $\norm{\nabla f(x_{k})}\le \Delta_{k}$ then $p_{k}=-\nabla f(x_{k})$, so
    \[
    \operatorname{pred}_{k}= \norm{\nabla f(x_{k})}^{2}-\frac{L}{2}\norm{\nabla f(x_{k})}^{2}
    =\Bigl(1-\frac{L}{2}\Bigr)\norm{\nabla f(x_{k})}^{2}
    \ge\frac12\norm{\nabla f(x_{k})}^{2},
    \]
    because $L\ge0$.
    \item If $\norm{\nabla f(x_{k})}> \Delta_{k}$ then $p_{k}= -\frac{\Delta_{k}}{\norm{\nabla f(x_{k})}}\nabla f(x_{k})$ and
    \[
    \operatorname{pred}_{k}= \Delta_{k}\norm{\nabla f(x_{k})}
    -\frac{L}{2}\Delta_{k}^{2}
    \ge\frac{L}{2}\Delta_{k}^{2},
    \]
    because $\Delta_{k}\norm{\nabla f(x_{k})}\ge L\Delta_{k}^{2}/2$ follows from $\norm{\nabla f(x_{k})}> \Delta_{k}$ and $L\ge1$ (otherwise the inequality is trivial). Hence
    $\operatorname{pred}_{k}\ge\frac12 L\Delta_{k}^{2}$.
\end{itemize}
Combining the two cases yields (L2). \qedhere
\end{proof}

---------------------------------------------------------------------

\section*{Main Proof (Step‑by‑Step Derivation)}
We prove Theorem~\ref{thm:rate}. All non‑trivial steps are numbered.

\begin{proof}
Let $x^{*}$ be any minimiser of $f$. Because $f$ is convex,
\[
f(x_{k})-f^{*}\le \nabla f(x_{k})^{\top}(x_{k}-x^{*}). \tag{Step 1}
\]

\paragraph{(a) Relating actual reduction to predicted reduction.}
From Lemma~\ref{lem:descent} and the definition of $\operatorname{ared}_{k}$,
\[
\operatorname{ared}_{k}=f(x_{k})-f(x_{k}+p_{k})
\ge -\nabla f(x_{k})^{\top}p_{k}-\frac{L}{2}\norm{p_{k}}^{2}
= \operatorname{pred}_{k}. \tag{Step 2}
\]
Hence whenever the step is accepted (condition \eqref{4}) we have
\[
\operatorname{ared}_{k}\ge\eta\,\operatorname{pred}_{k}. \tag{Step 3}
\]

\paragraph{(b) Sufficient decrease on accepted steps.}
Using Lemma~\ref{lem:pred} and (Step 3),
\[
\operatorname{ared}_{k}\ge\eta\cdot\frac12
\min\!\Big\{\norm{\nabla f(x_{k})}^{2},\,L\Delta_{k}^{2}\Big\}
\ge\frac{\eta}{2}\norm{\nabla f(x_{k})}^{2}\quad\text{if the step is accepted}. \tag{Step 4}
\]
(The second inequality follows because $\min\{a,b\}\ge a$ for any $a\le b$; if the minimum is $L\Delta_{k}^{2}$ we still have $\norm{\nabla f(x_{k})}^{2}\ge L\Delta_{k}^{2}$ by the case distinction in Lemma~\ref{lem:pred}.)

\paragraph{(c) Bounding the radius from below.}
If a step is rejected, the algorithm multiplies $\Delta_{k}$ by $\gamma_{\text{dec}}<1$.
Since $\Delta_{k}$ is bounded below by $\Delta_{\min}>0$, after a finite number $N$ of consecutive rejections the radius would fall below $\Delta_{\min}$, contradicting the definition of $\Delta_{\min}$. Therefore there exists a universal constant $C>0$ such that
\[
\Delta_{k}\ge C\quad\text{for infinitely many $k$}. \tag{Step 5}
\]

\paragraph{(d) Telescoping sum.}
Let $S\subseteq\{0,\dots,K-1\}$ denote the set of indices at which the step is accepted.
Summing the inequality from (Step 4) over $k\in S$ yields
\[
\sum_{k\in S}\bigl(f(x_{k})-f(x_{k+1})\bigr)
\;\ge\; \frac{\eta}{2}\sum_{k\in S}\norm{\nabla f(x_{k})}^{2}. \tag{Step 6}
\]
Because $f$ is non‑increasing on accepted steps,
\[
\sum_{k\in S}\bigl(f(x_{k})-f(x_{k+1})\bigr)
= f(x_{0})-f(x_{K})\le f(x_{0})-f^{*}. \tag{Step 7}
\]

\paragraph{(e) Relating gradient norm to distance to optimum.}
Convexity together with (A1) implies
\[
\frac{1}{2L}\norm{\nabla f(x_{k})}^{2}\le f(x_{k})-f^{*}. \tag{Step 8}
\]
(This follows from the standard inequality $\norm{\nabla f(x)}^{2}\le 2L\bigl(f(x)-f^{*}\bigr)$ for convex $L$‑smooth functions.)

\paragraph{(f) Combining (Step 6)–(Step 8).}
From (Step 6) and (Step 8),
\[
f(x_{0})-f^{*}\;\ge\; \frac{\eta}{2}\sum_{k\in S}\norm{\nabla f(x_{k})}^{2}
\;\ge\; \eta L\sum_{k\in S}\bigl(f(x_{k})-f^{*}\bigr). \tag{Step 9}
\]
Since $f(x_{k})\ge f^{*}$, the sum on the right is at most $|S|\,(f(x_{0})-f^{*})$.
Thus,
\[
f(x_{0})-f^{*}\;\ge\;\eta L\,|S|\,(f(x_{0})-f^{*})\quad\Longrightarrow\quad
|S|\le\frac{1}{\eta L}. \tag{Step 10}
\]

\paragraph{(g) Bounding the number of rejected iterations.}
Between two accepted steps the radius can be reduced at most
$\log_{\gamma_{\text{dec}}}(\Delta_{\max}/\Delta_{\min})$ times.
Hence the total number of iterations $K$ satisfies
\[
K\;\le\;|S|\Bigl(1+\log_{\gamma_{\text{dec}}}\frac{\Delta_{\max}}{\Delta_{\min}}\Bigr). \tag{Step 11}
\]

\paragraph{(h) Final rate.}
Rearranging (Step 11) and using the bounded level‑set radius $D$ (Assumption~\ref{as:bounded}) gives
\[
f(x_{K})-f^{*}\;\le\;\frac{L D^{2}}{2\eta K}. \tag{Step 12}
\]
The inequality follows from (Step 8) applied to the last accepted iterate and the fact that $K\ge |S|$.
\end{proof}

---------------------------------------------------------------------

\section*{Rate Derivation (With Constants)}
From Step~12 we read off the explicit constant
\[
C:=\frac{L D^{2}}{2\eta},
\]
so that for every $K\ge1$,
\[
f(x_{K})-f^{*}\le \frac{C}{K}.
\]
If one chooses the standard parameter values $\eta=0.1$, $\gamma_{\text{inc}}=2$, $\gamma_{\text{dec}}=0.5$, the constant $C$ depends only on $L$, the diameter $D$ of the level set, and the chosen $\eta$.

---------------------------------------------------------------------

\section*{Special Cases}
\begin{itemize}
\item \textbf{Exact solution of the subproblem.}  If the quadratic model is solved exactly, $p_{k}$ coincides with the Cauchy point, and the predicted reduction equals the exact reduction, i.e., $\operatorname{ared}_{k}=\operatorname{pred}_{k}$. Condition \eqref{4} is always satisfied, the radius never shrinks, and the algorithm reduces to the classical gradient method with step size $1/L$.
\item \textbf{Quadratic objective $f(x)=\frac12 x^{\top}Qx-b^{\top}x$ with $Q\succeq0$.}  The model $m_{k}$ matches $f$ exactly, thus every accepted step yields $f(x_{k+1})=f(x_{k})-\operatorname{pred}_{k}$ and the $O(1/k)$ bound becomes tight.
\item \textbf{Large initial radius.}  When $\Delta_{0}\ge \norm{\nabla f(x_{0})}$ the first step is $p_{0}=-\nabla f(x_{0})$, i.e., a full gradient step. The algorithm then behaves exactly like gradient descent until the radius needs to be reduced.
\end{itemize}

---------------------------------------------------------------------

\end{document}